%=========================
% Plantilla LaTeX Artículo Académico
% Compilar con: pdflatex + biber + pdflatex + pdflatex
%=========================
\documentclass[11pt,a4paper]{article}

%-------------------------
% Idioma y tipografía
%-------------------------
\usepackage[english]{babel}
\usepackage{csquotes}
\usepackage[T1]{fontenc}
\usepackage{lmodern}

%-------------------------
% Márgenes y espaciado
%-------------------------
\usepackage[a4paper,margin=2.5cm]{geometry}
\usepackage{setspace}
\onehalfspacing

%-------------------------
% Matemáticas y símbolos
%-------------------------
\usepackage{amsmath,amssymb,amsthm}

%-------------------------
% Tablas y figuras
%-------------------------
\usepackage{graphicx}
\usepackage{booktabs}
\usepackage{caption}
\usepackage{subcaption}
\usepackage{changepage}

%-------------------------
% Enlaces y PDF
%-------------------------
\usepackage{hyperref}
\hypersetup{
  colorlinks=true,
  linkcolor=blue,
  citecolor=blue,
  urlcolor=blue,
  pdftitle={The Continuum Hypothesis as a speech act},
  pdfauthor={Rodrigo Blanco Betes}
}

%-------------------------
% Listas, código y utilidades
%-------------------------
\usepackage{enumitem}
\usepackage{xcolor}
\usepackage{csquotes} % recomendado con biblatex

%-------------------------
% Bibliografía (biblatex + biber)
%-------------------------
\usepackage[
  backend=biber,
  style=apa,
  sortcites=true,
  url=true,
  doi=true
]{biblatex}
\DeclareLanguageMapping{english}{english-apa}
\addbibresource{refs.bib}

%-------------------------
% Metadatos del artículo
%-------------------------
\title{\textbf{The Continuum Hypothesis as a speech act}}

\author{Rodrigo Blanco Betes}


\date{\today}

%-------------------------
% Entornos de teoremas (opcional)
%-------------------------
\newtheorem{theorem}{Theorem}
\newtheorem{lemma}{Lemma}

\begin{document}
\maketitle

%-------------------------
% Resumen y palabras clave
%-------------------------
\begin{abstract}
A recent perspective by \textcite{Feferman2011a} has questioned whether continuum hypothesis (CH) is a well-defined and legitimate mathematical problem. This perspective has received recent criticism by P. Koellner (\citeyear{Koellner2016} \& \citeyear{Koellner2022}) for presenting the risk of reducing mathematics and science to the scope of subjectivism and convention. 

In this text, we approach this debate from a linguistic point of view, drawing on pragmatic theories of meaning and their application to mathematics. We focus particularly on speech act theory, and on the hypothesis—defended by Feferman—that mathematical objects function as elements of social reality, in the sense \textcite{searle1995} applies this concept to domains like politics and interpersonal relations. The purpose of this text is to address whether this point is justified, analyzing if speech act theory and Searle’s account of institutions could back up Feferman’s perspective, or at least contribute to a self-consistent account of mathematical reality. 
\end{abstract}

\noindent\textbf{Key words:} Continuum Hypothesis, Conceptual Structuralism, Speech Act Theory, Institution Theory.

%=========================
\section{Introduction}

In 1883, Cantor, after extending the concept of cardinality to infinite sets, posed the known as Continuum Hypothesis, stating that there could be no strictly intermediate set between the continuum and the set of the integers. After the development of axiomatic set theory, it was shown to be neither refutable (Gödel, \citeyear{Godel1940}) nor provable (Cohen, \citeyear{Cohen1966}) by the axioms of ZFC set theory, the ones more widely accepted. Later there have been many ways to try to fix the problem, but neither of them has been successful.

The logical and philosopher of mathematics Solomon Feferman proposed a new and controversial way to look at the problem. He developed a philosophical conception of mathematics called conceptual structuralism, which is based on the idea that mathematical notions originate in processes and activities from real experience that help us understand and model our environment. Around certain considerations, he was led to conclude that CH is an intrinsically vague problem, lacking any mathematical meaning.

This view has not been without its detractors. For example, views such as the ones of \textcite{Woodin2001a} or \textcite{Maddy1997} would directly contradict it, for considering the continuum problem to have a definite truth value, and the main goal of set-theorists to settle it through the introduction of new axioms. But Koellner (\citeyear{Koellner2016} \& \citeyear{Koellner2022}) expressed his disagreement in a much more direct way. One of the reasons for this disagreement was an important and peculiar feature of Feferman’s conception: its institutional account of the reality and objectivity of mathematical entities. Feferman, in the frame of conceptual structuralism, considers in (\citeyear{Feferman2011a}, p. 10) that mathematical entities’ objectivity lies upon their consistent appearance in independently working individuals; i.e. upon an “intersubjective objectivity that is ubiquitous in social reality”. He considers that mathematical entities exist objectively but are dependent on human constructions, following on this regard Searle (\citeyear{searle1995}, p.1): 

\begin{adjustwidth}{1cm}{1cm}
\small
\begin{quote}
[T]here are portions of the real world, objective facts in the world, that are only facts by human agreement. In a sense there are things that exist only because we believe them to exist (…). Things like money, property, governments, and marriages. Yet many facts regarding these things are ‘objective’ facts in the sense that they are not a matter of preferences, evaluations, or moral attitudes. 
\end{quote}
\end{adjustwidth}

In any case, he did not explain explicitly the way that mathematical entities could be socially constructed. While in Searle’s account we see quite explicitly, both in (\citeyear{searle1995}) and in (\citeyear{searle2010}), the development of a social institution theory in terms of speech act theory, Feferman simply points out that through a certain social conception of mathematics the truth value of determinate mathematical propositions, including CH, could be indeterminate: 

\begin{adjustwidth}{1cm}{1cm}
\small
\begin{quote}
My claim is that the basic conceptions of mathematics and their elaboration are also social constructions and that the objective reality that we ascribe to mathematics is simply the result of intersubjective objectivity (…). That is, it may simply be undecided under a given conception whether a given statement in the language of that conception has a determinate truth value. 
\end{quote}
\hfill \textcite[p. 13]{Feferman2011a}
\end{adjustwidth}

Maybe the lack of explicitness in this conception of mathematics is what led Koellner to criticize it. In his view, there is a huge difference between how political and social institutions come to exist and how numbers or mathematical entities do. For example, he defends (\citeyear{Koellner2022}, p. …) that the reality of the fact that “X is a Spanish citizen” is dependent on institutions such as governments, while the one of “there are infinitely many primes” is not. In general, he considers a huge mistake to make mathematical entities dependent on the intuition or the agreements of a determinate group of people, and that fully undermines, in his view, Feferman’s case on CH. 

My main thesis in this text will be that mathematical entities do share, as Feferman points out, many of the features of socially construed entities. In order to defend that, I am going to try to state in precise terms, following in most aspects Searle’s theory of social institutions, how the construction of mathematical entities (and, more precisely, domains) can be understood in a similar way to the establishment of social institutions.3 Once this parallelism has been properly addressed, it will follow in a natural way, I hold, that there are some previous preparatory conditions or presuppositions to the acceptance of a determinate mathematical entity. First, the entities introduced would have to be useful in a systematic way to solve, simplify, or generalize previously defined mathematical or scientific questions. For this reason, often the entities introduced should have a certain correspondence with a previously existent pre-theoretical concept, so that their application in those questions can be more straightforward. These conditions are somewhat coincident, and in any case compatible, with Feferman’s case that CH is an ill-defined mathematical problem, something I will explore in the end of the article. In general terms, I think that the institutional account for mathematical reality that I am going to explore in this text fits well with many aspects of Feferman’s conception of mathematics. 

\textbf{TEXT ORGANIZATION}

%=========================
\section{An institutional account for mathematical structures}

The controversy on the possibility that mathematical entities could have a social or institutional nature is not reduced to conceptual structuralism at all. In fact, there have been many accounts trying to justify that claim (Mention Wittgenstein, Lakatos, Paul Ernest). In recent times, one of the most impactful and explicit of these accounts has been the one of \textcite{Cole2013, Cole2008MathematicalDomains}, who, similarly to Feferman, follows the main Searle's notions on  social institutions. 

Searle's (\citeyear{searle1995} \& \citeyear{searle2010}) central argument is that we create social reality by collectively imposing status functions. According to this notion, institutional facets of reality are developed through a logical process that can be resumed by the formula $X$ counts as $Y$ in C, where $X$ is a physical object, $Y$ is a status function we assign to $X$ and C is the relevant context or community. This status function carries a set of deontic powers —rights, duties, obligations and authorizations. For example, the collective acceptance that a piece of paper $X$ counts as money $Y$ creates the right for its holder to make payments with it, to pay off a debt, to pay taxes, etc.  

\textcite{Cole2013} argues that Searle's framework of status functions is insufficient for a complete theory of institutions. To account for entities like borders, games, and literary characters, Cole introduces the concept of a representational function. He defines this as an institutional function whose primary purpose is to "facilitate our abilities to represent, analyze, reason about, [and] discover truths concerning... facets of reality that are not the entities in question" (p. 14). While status functions create deontic powers, representational functions equip us with cognitive tools to model and understand specific aspects of reality. For Cole, representational functions conform to a necessary component for explaining a vast range of institutional phenomena.\footnote{See \textcite{Cole2013} for more details and some examples.}

This distinction is crucial for our topic. The position of \textcite[1]{Cole2013} is that mathematical structures themselves are “freestanding institutional entities introduced to serve representational functions”. Although his literal phrasing often refers to “mathematical domains” he explicitly clarifies that his underlying commitment is to structures, stating that it is “mathematical structures that exist in virtue of collective agreement”.\footnote{One can see on this regard footnote \textcite[1]{Cole2013}: “More accurately, it is mathematical structures that exist in virtue of collective agreement. Yet, to avoid discussing structuralism, I shall write about mathematical domains”.} I understand such a structure as a specific arrangement of objects and the relations between them, which amounts to endowing a set with additional features—such as operations, topologies, or order relations. The term \textit{freestanding} refers to the fact that these representational functions are created \textit{ex nihilo}, without the need for a preexistent physical object these representational functions are assigned to (the $X$ in the previous formula). In summary, within Cole's framework, mathematical structures come to exist by being institutionally introduced in the form of representational functions, in such a way that the specific objects that appear in such structures —such as specific set theoretical representations of numbers— do not matter. This conception fits particularly well with the \textit{ante rem} structuralist view of mathematics, for which the specific nature of mathematical objects is non-essential; what is fundamental is the network of relations that constitutes the full structure. 

In Searle's framework, the idea of a social institution is closely related to a form of speech acts called Declarations, which have the unique property of changing the reality to match the propositional content of the speech act simply by successfully performing the act. The logical core of the formula X counts as Y in C is a Declaration. However, the imposition of a status or representational function is not always the result of a literal and explicit speech act. The most evident example is in language: it is an institution that provides both status and representational functions but cannot have been introduced by a set of explicit declarations, for those would presuppose language itself. In fact, Searle (...) considers language as the fundamental institution, which makes it possible that all other institutions come to exist.  

Sometimes we impose SFs and RFs using actual Declarations. The Captain of a soccer team might, for instance, `Declare Tom, you will be Our goalkeeper'. Yet this is not the only way in which Tom can become a Team's goalkeeper; he might simply stand by the relevant goal and, thereby, Cause all relevant parties to start talking and thinking about him as the Goalkeeper. Similarly, occasionally people formally Declare that someone Is their friend, yet typically they just start talking and thinking about that Person differently. As \textcite[10]{searle2010} puts it, on many occasions, `we Just linguistically treat or describe, or refer to, or talk about, or even think About an object in a way that creates a reality by representing that reality as Created'. Henceforth, ``Declaration'' shall refer to all such representational acts rather than merely Declarations in the aforementioned narrow sense. 

In other cases, institutions come to exist by tacit collective acceptance, habitual practice, or an informal agreement; these processes are what \textcite[...]{Cole2013} conceptualizes as standing Declarations. The underlying logical form of the resulting institutional fact remains that of a Declaration, even if it was not initiated by an explicit speech act. This point is radically important for mathematics. The most natural view is that mathematical axioms act as declarations that constitute mathematical reality, a perspective we see in \textcite[...]{SchmidtVenturi2024}. However, Gödel's incompleteness theorems demonstrate that there is no recursive set of first-order axioms that can determine up to isomorphism the most fundamental mathematical structures, such as the natural numbers or the real line.5 Therefore, if one accepts that mathematical structures are institutions in the sense defended here—and that most mathematicians work within the framework of first-order logic—then one must conclude that these structures are established and maintained not by a fixed set of axioms, but by standing Declarations of the mathematical community. 


"The number 3 is part of a system of symbols that we use to create and compute numbers. The existence of the number 3 is dependent on that system in a way that the existence of the moon is not." 

Coherence to mathematical practice: mathematical structures: N, R, C... 

Group theory 


%============================
\section{Preparatory conditions: is there a notion of \textit{mathematical authority}?}

The role of speech acts in modern mathematical axiomatics has been recently examined by \textcite{SchmidtVenturi2024}. A central thesis they defend is that a set of axioms serves as a declarative function. On this view, axioms are understood as bringing mathematical objects into existence—that is, as defining them. These axiomatic collections thus play a constitutive role, establishing the framework for a theory by acting as permissives that govern what can be legitimately asserted within a specific model. However, as I have argued in Section 2, mathematical structures are, in general, established prior to any specific set of axioms they may model. The intended structure—such as that of the natural numbers—guides the search for adequate axiomatizations, not the other way around. Therefore, it is not the axioms themselves, but a more fundamental standing Declaration—the collective, ongoing practice of mathematicians to identify and work within that structure—that ultimately defines the mathematical objects in question. In a similar vein, \textcite{RuffinoEtAl2021} offer one of the most systematic recent analyses of speech acts in mathematics. A central focus of their work is on mathematical declaratives, particularly definitions. For instance, they identify (p. 10084) consistency with prior definitions as a preparatory condition for successfully introducing a new one. Their analysis extends to many other dimensions of mathematical speech acts as they were defined by \textcite{SearleVanderveken1985}. 

However, these texts—and the literature at large—do not adequately analyze the preparatory conditions for axioms in their constitutive role, nor do they provide a clear account of such conditions for the standing Declarations that establish fundamental mathematical structures. This gap is crucial to address, as preparatory conditions are often the most determinative dimension of the declarative acts that ground our institutions. For example, most institutional declarations require the speaker to hold a special, recognized authority or to follow a specific procedure in the right circumstances. It seems immediately strange, however, to demand similar requirements for the introduction of a mathematical structure. The preparatory conditions for a standing Declaration that successfully institutes a new mathematical reality cannot be social or hierarchical in the usual sense. Yet, could there be a different sense of mathematical authority that serves as a preparatory condition for a felicitous institution of a mathematical structure?  

In general, authority is not an inherent quality but a conferred status; it arises when an individual's acts or pronouncements are consistently accepted as legitimate by a community, thereby constituting an institution. Furthermore, institutional authority is not just based on individual agents but also on the very procedures, roles, and instruments they employ. The legitimacy of an outcome often depends on its generation through an authoritative process. For example, a verdict delivered by a proper trial has much more authority than it would have if decided by a coin toss. Similarly, a proof in mathematics according to the rules of logic has authority, much more than a conclusion reached by a guess, even if it is compelling, or if an eminent mathematician has made it. In this sense, the means themselves can be seen as institutional agents, operating according to standing declarations that grant them specific powers and command collective acceptance. However, applying this framework immediately raises a deeper question: what specific criteria does legitimize the introduction of a new mathematical structure? There is no unanimously accepted view on what constitutes a legitimate, or felicitous, institutional declaration in mathematics. 

A clue can be given by the highly influential mathematician and philosopher of science Henri Poincaré, who can be considered a historical precursor to the institutional view of mathematics. Poincaré argued that mathematical entities are introduced by conventions. However, these are not arbitrary stipulations; they are guided by a principle of convenience and utility, making them fruitful for scientific and mathematical reasoning. He captures this idea eloquently: 

\begin{adjustwidth}{1cm}{1cm}
\begin{quote}
\small
Conventions are found above all in mathematics... It is precisely from this that they draw their rigor; our mind can affirm because it decrees. Are these decrees arbitrary? No, for if they were, they would be sterile. Our decrees are thus like those of an absolute, but wise, prince who would consult his Council of State. \parencite[3]{Poincare1952}
\end{quote}
\end{adjustwidth}

Following the terminology we are developing, Poincaré's decrees are institutional Declarations, while the “Council of State” represents the corresponding notion mathematical authority, which is a quite elucidatory metaphor; it represents the preparatory conditions of the Declarations introducing mathematical structures. In Poincaré's view (\citeyear{Poincare1952}, Chap. I), the fundamental convention of mathematics—one that is neither derivable from experience nor reducible to pure logic—is the acceptance of the principle of induction in arithmetic. This is essentially accepting that there is a foundational institutional Declaration that constitutes the very structure of the natural numbers. The “wisdom” of accepting this convention lies in its unparalleled fruitfulness for science as a whole. As Poincaré argues, deductive and analytical procedures merely rearrange existing knowledge, while induction is a unique engine of genuine discovery: 

\begin{adjustwidth}{1cm}{1cm}
\begin{quote}
\small
The analytical procedure of construction does not force us to descend, but it leaves us at the same level. We can only ascend by mathematical induction, which alone can teach us something new. Without the aid of this induction (...), construction would be powerless to create science. \parencite[28]{Poincare1952}
\end{quote}
\end{adjustwidth}

Poincaré's conventionalism had clear boundaries. He was profoundly reluctant to admit highly abstract entities and structures, such as those found in Cantorian set theory. His position made him a precursor of predicativism, arguing that only mathematical objects that can be predicatively defined—constructed step-by-step from the natural numbers without circular, self-referencing definitions—are legitimate.6  

As the main tittle of \textit{La science et la hypothése} suggests, mathematicians work with hypotheses. These hypotheses, according to Poincaré, are \textit{conventions}. Thus, these conventions to function rightly as hypotheses, must attain to two questions, (1) Why they are true, (2) What are they useful for.

Fruitfulness in mathematics cannot always be seen \textit{a-priori}, for it may depend on the obtained results. For example, the complex plane was accepted after some formulae were solved or the fundamental theorem of algebra solved.

For this reason, whether the introduction of a layer of mathematical reality is accepted could have to be decided \textit{a-posteriori}, so the process of mathematical authority is dynamic and iterative.

%=========================
\section{How could this affect the status of CH}

%=========================
\section{Criticism to the institutional conception of mathematics}

%=========================
\section{Conclusion}
What is to expect then from CH regarding this account? I don't think this is a question to be answered right now, but we can certainly approach it. First, the structures where CH can be formulated in have to be introduced: for example, the domain of third order arithmetic. Second, there is an evaluation of how this structure has been introduced, where a sense of \textit{mathematical authority} comes into play. Then, a the mathematical audience may judge, and decide whether this introduction is metodologically acceptable or not. Thus, a preparatory condition for this declaration, the one where the domain of third-order arithmetic was introduced, could come into play. 

\printbibliography

\end{document}
